%%%%%%%%%%%%%%%%%%%%%%%%%%%%%%%% FIXED 
\documentclass[11pt]{article}
\usepackage{graphicx}
\pagestyle{empty}
\setlength{\parskip}{0.25\baselineskip}
\renewcommand{\title}[1]{{\noindent\large\bfseries#1\medskip\\}}
\renewcommand{\author}[2]{{\noindent #1 \medskip\\ \small #2 \medskip\\}}
\renewcommand\refname{}
\usepackage[letterpaper,margin=20mm]{geometry}
\usepackage{hyperref}
%%%%%%%%%%%%%%%%%%%%%%%%%%%%%%%%

\begin{document}

\title{Projecting possible future trajectories for systems of cities}
\author{
% Authors Names
Juste Raimbault,\textsuperscript{1,2,3,4}
Denise Pumain,\textsuperscript{4}
}
{
% Authors Affiliations
1. LaSTIG, Univ. Gustave Eiffel, IGN-ENSG\\
2. CASA, UCL\\
3. UPS CNRS 3611 ISC-PIF\\
4. UMR CNRS 8504 G{\'e}ographie-cit{\'e}s\\
}


% mini abstract 2000 caracs
% 

% ABSTRACT

Designing sustainable policies for territorial systems requires an approach encompassing both generic processes and geographical peculiarities of diverse urban systems across the world, in a multi-level way \cite{rozenblat2018conclusion},% TODO relire - pas sur que suggere simu models -> cite book chapter handbook entropy?
which can further be quantified through simulation models for systems of cities \cite{raimbault2021spatial}. An evolutionary theory of urban systems 

Methods for exploring various contingencies in such simulation models have been developed in close link to the validation of these models: the MARIUS model family coupling population growth with economic exchanges \cite{} has been challenged into ``unexpected'' behaviours with the Pattern Space Exploration algorithm \cite{cherel2015beyond}. Similarly, the SimpopNet model for the co-evolution of transportation infrastructure networks and cities  ... \cite{raimbault2020unveiling}

We build in this contribution on these previously developed models and exploration methods \cite{raimbault2019methods}, to introduce a new composite method for projecting possible future plausible trajectories for a system of cities. 
\cite{raimbault2020indirect} %showed that a set of parameter points

% We first calibrate the full macroscopic model (all submodels included) on population trajectories, providing a range of possible parametrisations having produced the past dynamics, with the points on the Pareto front (bi-objective calibration on log(MSE) - MSE(log);
%From these optimal points, extract parameter boundaries and consider these as the most probable values;
%Use these parameter bounds to run the PSE diversity search algorithms, on dynamics a few time steps in the future to obtain the universe of possibilities, what quantifies possible regimes (in the case of a high dimensional output space or a high number of trajectories, it is interpreted as adaptability [?] - a restricted future implies less adaptability)

Our method is applied to the European and Chinese system of cities, using harmonised datasets for population trajectories of cities constructed by \cite{cura2017old}, and in particular the ChinaCities database \cite{swerts2017data}.

Simulation results are shown in Figure 1.
% TODO find a specific representation, with incertitude tubes, for population, and other SDG dimensions? or radar plot, how the system evolves on it? or most important dim only and project? -> emissions, GDP? ~ ! possible that with such a high dimensionality the model can tell anything - finding trajs with PSE (then how to be sure to remove bias in visu or interpretation?)



\newpage


\begin{figure}%[b]
  \centering
  %\includegraphics[width=\textwidth]{figures/}
  \caption{}
\end{figure}

{\small
%\linespread{1}\selectfont
\bibliographystyle{unsrt}
\bibliography{biblio}
}

\end{document}




%\bigskip
%{\small
%\noindent[1] Here you can add references.\\
%\noindent[2] Differently from this one, not written using a LLM.
%}

%\begin{figure}[b]
%  \centering
%  \includegraphics[width=0.5\textwidth]{figure.pdf}
%  \caption{Figure caption.}
%\end{figure}

